\chapter{Unit 6: Multiple Integrals and Vector Volume Calculus}

\section{Section A: Multiple Choice Practice Questions (60 Curated Problems)}

\noindent\textbf{Q1. What does a double integral $\iint$ R f ( x , y ) d x d y represent geometrically if f ( x , y ) = 1 ?} \hfill [Tag: double integrals | Easy]
\begin{enumerate}[label=(\alph*),topsep=2pt,itemsep=1pt,partopsep=0pt,parsep=0pt]
    \item {{The length of the boundary of R}}
    \item {{The volume under the surface f ( x , y )}}
    \item {{The centroid of the region R}}
    \item {{The area of the region R}}
\end{enumerate}
\vspace{2mm}

\noindent\textbf{Q2. In a double integral $\int$ a b $\int$ c d f ( x , y ) d y d x , which variable is integrated first?} \hfill [Tag: double integrals | Easy]
\begin{enumerate}[label=(\alph*),topsep=2pt,itemsep=1pt,partopsep=0pt,parsep=0pt]
    \item {{x}}
    \item {{Either x}}
    \item {{Neither}}
    \item {{y}}
\end{enumerate}
\vspace{2mm}

\noindent\textbf{Q3. Evaluate the double integral $\int$ 0 1 $\int$ 0 2 1 d y d x .} \hfill [Tag: double integrals | Easy]
\begin{enumerate}[label=(\alph*),topsep=2pt,itemsep=1pt,partopsep=0pt,parsep=0pt]
    \item {{4}}
    \item {{2}}
    \item {{1}}
    \item {{3}}
\end{enumerate}
\vspace{2mm}

\noindent\textbf{Q4. What is the primary purpose of changing the order of integration in a double integral?} \hfill [Tag: change of order of integration | Easy]
\begin{enumerate}[label=(\alph*),topsep=2pt,itemsep=1pt,partopsep=0pt,parsep=0pt]
    \item {{To change a double integral into a triple integral}}
    \item {{To change the value of the integral}}
    \item {{To simplify the evaluation of the integral}}
    \item {{To find the perimeter of a region}}
\end{enumerate}
\vspace{2mm}

\noindent\textbf{Q5. If the original limits of integration are 0 $\le$ x $\le$ 1 and 0 $\le$ y $\le$ x , what are the new limits if the order is changed to d x d y ?} \hfill [Tag: change of order of integration | Easy]
\begin{enumerate}[label=(\alph*),topsep=2pt,itemsep=1pt,partopsep=0pt,parsep=0pt]
    \item {{0 $\le$}}
    \item {{0 $\le$}}
    \item {{0 $\le$}}
    \item {{0 $\le$}}
\end{enumerate}
\vspace{2mm}

\noindent\textbf{Q6. Does changing the order of integration alter the final value of a definite double integral over a closed region?} \hfill [Tag: change of order of integration | Easy]
\begin{enumerate}[label=(\alph*),topsep=2pt,itemsep=1pt,partopsep=0pt,parsep=0pt]
    \item {{No, the value remains exactly the same.}}
    \item {{Yes, it produces a different numerical value.}}
    \item {{Only if the integrand contains trigonometric functions.}}
    \item {{Yes, it always changes the sign.}}
\end{enumerate}
\vspace{2mm}

\noindent\textbf{Q7. When converting a double integral from Cartesian ( x , y ) to polar ( r , $\theta$ ) coordinates, what does d x d y become?} \hfill [Tag: change of variables | Easy]
\begin{enumerate}[label=(\alph*),topsep=2pt,itemsep=1pt,partopsep=0pt,parsep=0pt]
    \item {{r 2}}
    \item {{r}}
    \item {{d r d $\theta$}}
    \item {{r d r d $\theta$}}
\end{enumerate}
\vspace{2mm}

\noindent\textbf{Q8. What is the Jacobian J = $\partial$ ( u , v ) $\partial$ ( x , y ) used for in double integrals?} \hfill [Tag: change of variables | Easy]
\begin{enumerate}[label=(\alph*),topsep=2pt,itemsep=1pt,partopsep=0pt,parsep=0pt]
    \item {{To find the limits of integration}}
    \item {{To convert a double integral into a single integral}}
    \item {{To scale the differential area element when changing variables}}
    \item {{To integrate polynomials}}
\end{enumerate}
\vspace{2mm}

\noindent\textbf{Q9. In polar coordinates, what is the relation for x 2 + y 2 ?} \hfill [Tag: change of variables | Easy]
\begin{enumerate}[label=(\alph*),topsep=2pt,itemsep=1pt,partopsep=0pt,parsep=0pt]
    \item {{r cos $\theta$}}
    \item {{$\theta$}}
    \item {{r}}
    \item {{r 2}}
\end{enumerate}
\vspace{2mm}

\noindent\textbf{Q10. Which of the following formulas calculates the area of a region R in the x y -plane?} \hfill [Tag: application of double integrals to calculate area and volume | Easy]
\begin{enumerate}[label=(\alph*),topsep=2pt,itemsep=1pt,partopsep=0pt,parsep=0pt]
    \item {{$\iint$ R}}
    \item {{$\iint$ R}}
    \item {{$\iint$ R}}
    \item {{$\iint$ R}}
\end{enumerate}
\vspace{2mm}

\noindent\textbf{Q11. If f ( x , y ) > 0 over a region R , what does $\iint$ R f ( x , y ) d A compute?} \hfill [Tag: application of double integrals to calculate area and volume | Easy]
\begin{enumerate}[label=(\alph*),topsep=2pt,itemsep=1pt,partopsep=0pt,parsep=0pt]
    \item {{The area of the region R}}
    \item {{The volume of the solid bounded above by z =}}
    \item {{The perimeter of the region R}}
    \item {{The surface area of f ( x , y )}}
\end{enumerate}
\vspace{2mm}

\noindent\textbf{Q12. A triple integral is typically used to integrate a function over a region in how many dimensions?} \hfill [Tag: triple integrals | Easy]
\begin{enumerate}[label=(\alph*),topsep=2pt,itemsep=1pt,partopsep=0pt,parsep=0pt]
    \item {{Three}}
    \item {{Four}}
    \item {{Two}}
    \item {{One}}
\end{enumerate}
\vspace{2mm}

\noindent\textbf{Q13. What is the standard differential volume element d V in Cartesian coordinates?} \hfill [Tag: triple integrals | Easy]
\begin{enumerate}[label=(\alph*),topsep=2pt,itemsep=1pt,partopsep=0pt,parsep=0pt]
    \item {{$\rho$ 2}}
    \item {{d r d $\theta$ d z}}
    \item {{d x d y}}
    \item {{d x d y d z}}
\end{enumerate}
\vspace{2mm}

\noindent\textbf{Q14. Evaluate $\int$ 0 1 $\int$ 0 1 $\int$ 0 1 1 d z d y d x .} \hfill [Tag: triple integrals | Easy]
\begin{enumerate}[label=(\alph*),topsep=2pt,itemsep=1pt,partopsep=0pt,parsep=0pt]
    \item {{1/3}}
    \item {{1}}
    \item {{3}}
    \item {{0}}
\end{enumerate}
\vspace{2mm}

\noindent\textbf{Q15. How can the volume of a 3D solid region D be found using a triple integral?} \hfill [Tag: application of triple integrals to calculate volume | Easy]
\begin{enumerate}[label=(\alph*),topsep=2pt,itemsep=1pt,partopsep=0pt,parsep=0pt]
    \item {{$\iiint$ D}}
    \item {{$\iiint$ D}}
    \item {{$\iiint$ D}}
    \item {{$\iiint$ D}}
\end{enumerate}
\vspace{2mm}

\noindent\textbf{Q16. When finding the volume of a sphere, which coordinate system is usually the most convenient to use for the triple integral?} \hfill [Tag: application of triple integrals to calculate volume | Easy]
\begin{enumerate}[label=(\alph*),topsep=2pt,itemsep=1pt,partopsep=0pt,parsep=0pt]
    \item {{Polar coordinates}}
    \item {{Spherical coordinates}}
    \item {{Cylindrical coordinates}}
    \item {{Cartesian coordinates}}
\end{enumerate}
\vspace{2mm}

\noindent\textbf{Q17. In calculating the area of a circle of radius R using polar coordinates, what are the limits for r and $\theta$ ?} \hfill [Tag: application of double integrals to calculate area and volume | Easy]
\begin{enumerate}[label=(\alph*),topsep=2pt,itemsep=1pt,partopsep=0pt,parsep=0pt]
    \item {{- R $\le$}}
    \item {{0 $\le$}}
    \item {{0 $\le$}}
    \item {{0 $\le$}}
\end{enumerate}
\vspace{2mm}

\noindent\textbf{Q18. What is the Jacobian determinant J for converting Cartesian coordinates to cylindrical coordinates ( r , $\theta$ , z ) ?} \hfill [Tag: change of variables | Easy]
\begin{enumerate}[label=(\alph*),topsep=2pt,itemsep=1pt,partopsep=0pt,parsep=0pt]
    \item {{r}}
    \item {{r 2}}
    \item {{r sin $\theta$}}
    \item {{1}}
\end{enumerate}
\vspace{2mm}

\noindent\textbf{Q19. What does the differential volume element d V become in spherical coordinates ( $\rho$ , $\phi$ , $\theta$ ) ?} \hfill [Tag: application of triple integrals to calculate volume | Easy]
\begin{enumerate}[label=(\alph*),topsep=2pt,itemsep=1pt,partopsep=0pt,parsep=0pt]
    \item {{r d r d $\theta$ d z}}
    \item {{$\rho$ 2}}
    \item {{$\rho$ sin $\phi$ d $\rho$ d $\phi$ d $\theta$}}
    \item {{$\rho$ 2}}
\end{enumerate}
\vspace{2mm}

\noindent\textbf{Q20. If a region R is a rectangle defined by 1 $\le$ x $\le$ 3 and 2 $\le$ y $\le$ 5 , what is the area of R computed by $\iint$ R d A ?} \hfill [Tag: double integrals | Easy]
\begin{enumerate}[label=(\alph*),topsep=2pt,itemsep=1pt,partopsep=0pt,parsep=0pt]
    \item {{15}}
    \item {{8}}
    \item {{5}}
    \item {{6}}
\end{enumerate}
\vspace{2mm}

\noindent\textbf{Q21. Evaluate the double integral $\int$ 0 1 $\int$ 0 x ( x 2 + y ) d y d x .} \hfill [Tag: double integrals | Medium]
\begin{enumerate}[label=(\alph*),topsep=2pt,itemsep=1pt,partopsep=0pt,parsep=0pt]
    \item {{3}}
    \item {{12}}
    \item {{4}}
    \item {{12}}
\end{enumerate}
\vspace{2mm}

\noindent\textbf{Q22. What is the value of $\int$ 0 2 $\int$ 0 1 ( 2 x + 3 y ) d x d y ?} \hfill [Tag: double integrals | Medium]
\begin{enumerate}[label=(\alph*),topsep=2pt,itemsep=1pt,partopsep=0pt,parsep=0pt]
    \item {{4}}
    \item {{8}}
    \item {{6}}
    \item {{10}}
\end{enumerate}
\vspace{2mm}

\noindent\textbf{Q23. Consider the integral $\int$ 0 1 $\int$ x 1 f ( x , y ) d y d x . If the order of integration is changed, what are the new limits?} \hfill [Tag: change of order of integration | Medium]
\begin{enumerate}[label=(\alph*),topsep=2pt,itemsep=1pt,partopsep=0pt,parsep=0pt]
    \item {{$\int$ 0}}
    \item {{$\int$ 0}}
    \item {{$\int$ 0}}
    \item {{$\int$ 0}}
\end{enumerate}
\vspace{2mm}

\noindent\textbf{Q24. When changing the order of integration for $\int$ 0 2 $\int$ y 2 /2 y f ( x , y ) d x d y , what is the outer integral limit for x ?} \hfill [Tag: change of order of integration | Medium]
\begin{enumerate}[label=(\alph*),topsep=2pt,itemsep=1pt,partopsep=0pt,parsep=0pt]
    \item {{0}}
    \item {{0}}
    \item {{0}}
    \item {{0}}
\end{enumerate}
\vspace{2mm}

\noindent\textbf{Q25. What is the Jacobian $\partial$ ( r , $\theta$ ) $\partial$ ( x , y ) for the transformation to polar coordinates x = r cos $\theta$ , y = r sin $\theta$ ?} \hfill [Tag: change of variables | Medium]
\begin{enumerate}[label=(\alph*),topsep=2pt,itemsep=1pt,partopsep=0pt,parsep=0pt]
    \item {{sin $\theta$ cos $\theta$}}
    \item {{1}}
    \item {{r}}
    \item {{r 2}}
\end{enumerate}
\vspace{2mm}

\noindent\textbf{Q26. Evaluate $\iint$ R e x 2 + y 2 d x d y where R is the region x 2 + y 2 $\le$ 4 in the first quadrant, using polar coordinates.} \hfill [Tag: change of variables | Medium]
\begin{enumerate}[label=(\alph*),topsep=2pt,itemsep=1pt,partopsep=0pt,parsep=0pt]
    \item {{$\pi$ ( e 4}}
    \item {{4}}
    \item {{2}}
    \item {{4}}
\end{enumerate}
\vspace{2mm}

\noindent\textbf{Q27. Use a double integral to find the area of the region bounded by y = x 2 and y = x .} \hfill [Tag: application of double integrals to calculate area and volume | Medium]
\begin{enumerate}[label=(\alph*),topsep=2pt,itemsep=1pt,partopsep=0pt,parsep=0pt]
    \item {{2}}
    \item {{6}}
    \item {{3}}
    \item {{4}}
\end{enumerate}
\vspace{2mm}

\noindent\textbf{Q28. Find the volume under the surface z = x y and above the rectangle 0 $\le$ x $\le$ 2 , 0 $\le$ y $\le$ 3 .} \hfill [Tag: application of double integrals to calculate area and volume | Medium]
\begin{enumerate}[label=(\alph*),topsep=2pt,itemsep=1pt,partopsep=0pt,parsep=0pt]
    \item {{9}}
    \item {{12}}
    \item {{18}}
    \item {{6}}
\end{enumerate}
\vspace{2mm}

\noindent\textbf{Q29. Evaluate the triple integral $\int$ 0 1 $\int$ 0 x $\int$ 0 y 2 z d z d y d x .} \hfill [Tag: triple integrals | Medium]
\begin{enumerate}[label=(\alph*),topsep=2pt,itemsep=1pt,partopsep=0pt,parsep=0pt]
    \item {{12}}
    \item {{8}}
    \item {{24}}
    \item {{6}}
\end{enumerate}
\vspace{2mm}

\noindent\textbf{Q30. What is the value of $\int$ 0 1 $\int$ 0 1 $\int$ 0 1 ( x + y + z ) d x d y d z ?} \hfill [Tag: triple integrals | Medium]
\begin{enumerate}[label=(\alph*),topsep=2pt,itemsep=1pt,partopsep=0pt,parsep=0pt]
    \item {{1}}
    \item {{2}}
    \item {{2}}
    \item {{2}}
\end{enumerate}
\vspace{2mm}

\noindent\textbf{Q31. Find the volume of the solid bounded by the planes x = 0 , y = 0 , z = 0 , and x + y + z = 1 .} \hfill [Tag: application of triple integrals to calculate volume | Medium]
\begin{enumerate}[label=(\alph*),topsep=2pt,itemsep=1pt,partopsep=0pt,parsep=0pt]
    \item {{3}}
    \item {{6}}
    \item {{4}}
    \item {{2}}
\end{enumerate}
\vspace{2mm}

\noindent\textbf{Q32. Calculate the volume of a sphere of radius R using triple integration in spherical coordinates.} \hfill [Tag: application of triple integrals to calculate volume | Medium]
\begin{enumerate}[label=(\alph*),topsep=2pt,itemsep=1pt,partopsep=0pt,parsep=0pt]
    \item {{3}}
    \item {{4 $\pi$ R 3}}
    \item {{3}}
    \item {{3}}
\end{enumerate}
\vspace{2mm}

\noindent\textbf{Q33. Evaluate $\int$ 0 $\pi$ /2 $\int$ 0 c o s x y sin x d y d x .} \hfill [Tag: double integrals | Medium]
\begin{enumerate}[label=(\alph*),topsep=2pt,itemsep=1pt,partopsep=0pt,parsep=0pt]
    \item {{4}}
    \item {{2}}
    \item {{6}}
    \item {{3}}
\end{enumerate}
\vspace{2mm}

\noindent\textbf{Q34. The integral $\int$ 0 4 $\int$ x 2 f ( x , y ) d y d x after changing the order of integration becomes:} \hfill [Tag: change of order of integration | Medium]
\begin{enumerate}[label=(\alph*),topsep=2pt,itemsep=1pt,partopsep=0pt,parsep=0pt]
    \item {{$\int$ 0}}
    \item {{$\int$ 0}}
    \item {{$\int$ 0}}
    \item {{$\int$ 0}}
\end{enumerate}
\vspace{2mm}

\noindent\textbf{Q35. Evaluate the integral $\iint$ D ( x 2 + y 2 ) 3/2 d A where D is the disk x 2 + y 2 $\le$ 1 using polar coordinates.} \hfill [Tag: change of variables | Medium]
\begin{enumerate}[label=(\alph*),topsep=2pt,itemsep=1pt,partopsep=0pt,parsep=0pt]
    \item {{5}}
    \item {{3}}
    \item {{$\pi$}}
    \item {{5}}
\end{enumerate}
\vspace{2mm}

\noindent\textbf{Q36. Find the area of the circle x 2 + y 2 = a 2 using double integration.} \hfill [Tag: application of double integrals to calculate area and volume | Medium]
\begin{enumerate}[label=(\alph*),topsep=2pt,itemsep=1pt,partopsep=0pt,parsep=0pt]
    \item {{3}}
    \item {{2 $\pi$a}}
    \item {{2}}
    \item {{$\pi$ a 2}}
\end{enumerate}
\vspace{2mm}

\noindent\textbf{Q37. Evaluate $\iiint$ E x d V where E is the region 0 $\le$ x $\le$ 1 , 0 $\le$ y $\le$ 2 , 0 $\le$ z $\le$ 3 .} \hfill [Tag: triple integrals | Medium]
\begin{enumerate}[label=(\alph*),topsep=2pt,itemsep=1pt,partopsep=0pt,parsep=0pt]
    \item {{2}}
    \item {{1.5}}
    \item {{6}}
    \item {{3}}
\end{enumerate}
\vspace{2mm}

\noindent\textbf{Q38. Calculate the volume of the cylinder x 2 + y 2 = 4 bounded by the planes z = 0 and z = 5 .} \hfill [Tag: application of triple integrals to calculate volume | Medium]
\begin{enumerate}[label=(\alph*),topsep=2pt,itemsep=1pt,partopsep=0pt,parsep=0pt]
    \item {{40 $\pi$}}
    \item {{10 $\pi$}}
    \item {{20 $\pi$}}
    \item {{5 $\pi$}}
\end{enumerate}
\vspace{2mm}

\noindent\textbf{Q39. If transforming coordinates using x = u + v and y = u - v , what is the absolute value of the Jacobian | J | ?} \hfill [Tag: change of variables | Medium]
\begin{enumerate}[label=(\alph*),topsep=2pt,itemsep=1pt,partopsep=0pt,parsep=0pt]
    \item {{- 2}}
    \item {{u 2}}
    \item {{2}}
    \item {{1}}
\end{enumerate}
\vspace{2mm}

\noindent\textbf{Q40. Evaluate $\int$ 0 1 $\int$ y 1 e x 2 d x d y .} \hfill [Tag: double integrals | Medium]
\begin{enumerate}[label=(\alph*),topsep=2pt,itemsep=1pt,partopsep=0pt,parsep=0pt]
    \item {{2}}
    \item {{e -}}
    \item {{2}}
    \item {{2}}
\end{enumerate}
\vspace{2mm}

\noindent\textbf{Q41. Evaluate the integral $\int$ 0 $\infty$ $\int$ 0 $\infty$ e - ( x 2 + y 2 ) cos ( x 2 + y 2 ) d x d y .} \hfill [Tag: double integrals | Hard]
\begin{enumerate}[label=(\alph*),topsep=2pt,itemsep=1pt,partopsep=0pt,parsep=0pt]
    \item {{8}}
    \item {{8}}
    \item {{4 2}}
    \item {{4}}
\end{enumerate}
\vspace{2mm}

\noindent\textbf{Q42. What is the value of the double integral $\iint$ D | x - y | d x d y , where D = {( x , y ) | 0 $\le$ x $\le$ 1 , 0 $\le$ y $\le$ 1 } ?} \hfill [Tag: double integrals | Hard]
\begin{enumerate}[label=(\alph*),topsep=2pt,itemsep=1pt,partopsep=0pt,parsep=0pt]
    \item {{3}}
    \item {{4}}
    \item {{6}}
    \item {{2}}
\end{enumerate}
\vspace{2mm}

\noindent\textbf{Q43. Evaluate lim n $\to$ $\infty$ $\iint$ D ( x 2 + y 2 ) n d x d y where D is the unit disk x 2 + y 2 $\le$ 1 .} \hfill [Tag: double integrals | Hard]
\begin{enumerate}[label=(\alph*),topsep=2pt,itemsep=1pt,partopsep=0pt,parsep=0pt]
    \item {{0}}
    \item {{$\pi$}}
    \item {{2}}
    \item {{$\infty$}}
\end{enumerate}
\vspace{2mm}

\noindent\textbf{Q44. Change the order of integration for $\int$ - 1 1 $\int$ x 2 1 f ( x , y ) d y d x . Which of the following is the correct equivalent integral?} \hfill [Tag: change of order of integration | Hard]
\begin{enumerate}[label=(\alph*),topsep=2pt,itemsep=1pt,partopsep=0pt,parsep=0pt]
    \item {{$\int$ - 1}}
    \item {{$\int$ 0}}
    \item {{$\int$ 0}}
    \item {{$\int$ - 1}}
\end{enumerate}
\vspace{2mm}

\noindent\textbf{Q45. Evaluate $\int$ 0 2 $\int$ y /2 1 e x 2 d x d y by reversing the order of integration.} \hfill [Tag: change of order of integration | Hard]
\begin{enumerate}[label=(\alph*),topsep=2pt,itemsep=1pt,partopsep=0pt,parsep=0pt]
    \item {{2 ( e -}}
    \item {{2}}
    \item {{e 2}}
    \item {{e -}}
\end{enumerate}
\vspace{2mm}

\noindent\textbf{Q46. Consider the integral $\int$ 0 1 $\int$ e x e f ( x , y ) d y d x . After changing the order of integration, what are the new limits?} \hfill [Tag: change of order of integration | Hard]
\begin{enumerate}[label=(\alph*),topsep=2pt,itemsep=1pt,partopsep=0pt,parsep=0pt]
    \item {{$\int$ 0}}
    \item {{$\int$ 1}}
    \item {{$\int$ 0}}
    \item {{$\int$ 1}}
\end{enumerate}
\vspace{2mm}

\noindent\textbf{Q47. If I = $\int$ 0 $\pi$ $\int$ x $\pi$ y s i n y d y d x , find the value of I using change of order.} \hfill [Tag: change of order of integration | Hard]
\begin{enumerate}[label=(\alph*),topsep=2pt,itemsep=1pt,partopsep=0pt,parsep=0pt]
    \item {{2}}
    \item {{$\pi$}}
    \item {{1}}
    \item {{0}}
\end{enumerate}
\vspace{2mm}

\noindent\textbf{Q48. Find the Jacobian J = $\partial$ ( u , v ) $\partial$ ( x , y ) for the transformation x = u 2 - v 2 and y = 2 uv .} \hfill [Tag: change of variables | Hard]
\begin{enumerate}[label=(\alph*),topsep=2pt,itemsep=1pt,partopsep=0pt,parsep=0pt]
    \item {{4 uv}}
    \item {{2 ( u 2}}
    \item {{u 2}}
    \item {{4 ( u 2}}
\end{enumerate}
\vspace{2mm}

\noindent\textbf{Q49. Use the transformation u = x + y , v = x - y to evaluate $\iint$ R ( x + y ) 2 e x 2 - y 2 d x d y , where R is the region bounded by x + y = 1 , x + y = 2 , x - y = - 1 , x - y = 1 .} \hfill [Tag: change of variables | Hard]
\begin{enumerate}[label=(\alph*),topsep=2pt,itemsep=1pt,partopsep=0pt,parsep=0pt]
    \item {{2}}
    \item {{4}}
    \item {{4}}
    \item {{2}}
\end{enumerate}
\vspace{2mm}

\noindent\textbf{Q50. Evaluate the integral $\iint$ R sin ( x + y x - y ) d x d y over the triangle R with vertices ( 0 , 0 ) , ( 1 , 0 ) , ( 0 , 1 ) using the substitution u = x - y , v = x + y .} \hfill [Tag: change of variables | Hard]
\begin{enumerate}[label=(\alph*),topsep=2pt,itemsep=1pt,partopsep=0pt,parsep=0pt]
    \item {{2}}
    \item {{2}}
    \item {{4}}
    \item {{4}}
\end{enumerate}
\vspace{2mm}

\noindent\textbf{Q51. Find the volume of the solid bounded by the paraboloid z = 4 - x 2 - y 2 and the x y -plane.} \hfill [Tag: application of double integrals to calculate area and volume | Hard]
\begin{enumerate}[label=(\alph*),topsep=2pt,itemsep=1pt,partopsep=0pt,parsep=0pt]
    \item {{32 $\pi$ /3}}
    \item {{4 $\pi$}}
    \item {{16 $\pi$}}
    \item {{8 $\pi$}}
\end{enumerate}
\vspace{2mm}

\noindent\textbf{Q52. Calculate the area of the region bounded by the lemniscate r 2 = a 2 cos 2 $\theta$ .} \hfill [Tag: application of double integrals to calculate area and volume | Hard]
\begin{enumerate}[label=(\alph*),topsep=2pt,itemsep=1pt,partopsep=0pt,parsep=0pt]
    \item {{2}}
    \item {{2 a 2}}
    \item {{$\pi$ a 2}}
    \item {{a 2}}
\end{enumerate}
\vspace{2mm}

\noindent\textbf{Q53. Find the volume common to the cylinders x 2 + y 2 = a 2 and x 2 + z 2 = a 2 .} \hfill [Tag: application of double integrals to calculate area and volume | Hard]
\begin{enumerate}[label=(\alph*),topsep=2pt,itemsep=1pt,partopsep=0pt,parsep=0pt]
    \item {{8 a 3}}
    \item {{3}}
    \item {{3}}
    \item {{3}}
\end{enumerate}
\vspace{2mm}

\noindent\textbf{Q54. Determine the area of the region outside the circle r = 2 and inside the cardioid r = 2 ( 1 + cos $\theta$ ) .} \hfill [Tag: application of double integrals to calculate area and volume | Hard]
\begin{enumerate}[label=(\alph*),topsep=2pt,itemsep=1pt,partopsep=0pt,parsep=0pt]
    \item {{8 +}}
    \item {{8 -}}
    \item {{2 +}}
    \item {{4 +}}
\end{enumerate}
\vspace{2mm}

\noindent\textbf{Q55. Evaluate $\iiint$ E z d V , where E is the solid tetrahedron bounded by the planes x = 0 , y = 0 , z = 0 , and x + y + z = 1 .} \hfill [Tag: triple integrals | Hard]
\begin{enumerate}[label=(\alph*),topsep=2pt,itemsep=1pt,partopsep=0pt,parsep=0pt]
    \item {{24}}
    \item {{48}}
    \item {{6}}
    \item {{12}}
\end{enumerate}
\vspace{2mm}

\noindent\textbf{Q56. Evaluate $\iiint$ V x 2 d V where V is the region bounded by x 2 + y 2 + z 2 = 1 .} \hfill [Tag: triple integrals | Hard]
\begin{enumerate}[label=(\alph*),topsep=2pt,itemsep=1pt,partopsep=0pt,parsep=0pt]
    \item {{5}}
    \item {{15}}
    \item {{15}}
    \item {{15}}
\end{enumerate}
\vspace{2mm}

\noindent\textbf{Q57. Using cylindrical coordinates, set up the triple integral for the volume of the solid bounded above by the sphere x 2 + y 2 + z 2 = 2 and below by the paraboloid z = x 2 + y 2 .} \hfill [Tag: triple integrals | Hard]
\begin{enumerate}[label=(\alph*),topsep=2pt,itemsep=1pt,partopsep=0pt,parsep=0pt]
    \item {{$\int$ 0}}
    \item {{$\int$ 0}}
    \item {{$\int$ 0}}
    \item {{$\int$ 0}}
\end{enumerate}
\vspace{2mm}

\noindent\textbf{Q58. Find the volume of the region bounded by the ellipsoid a 2 x 2 + b 2 y 2 + c 2 z 2 = 1 .} \hfill [Tag: application of triple integrals to calculate volume | Hard]
\begin{enumerate}[label=(\alph*),topsep=2pt,itemsep=1pt,partopsep=0pt,parsep=0pt]
    \item {{3}}
    \item {{3}}
    \item {{$\pi$ab c}}
    \item {{3}}
\end{enumerate}
\vspace{2mm}

\noindent\textbf{Q59. Calculate the volume of the solid bounded by the cone z = x 2 + y 2 and the sphere x 2 + y 2 + z 2 = z .} \hfill [Tag: application of triple integrals to calculate volume | Hard]
\begin{enumerate}[label=(\alph*),topsep=2pt,itemsep=1pt,partopsep=0pt,parsep=0pt]
    \item {{6}}
    \item {{4}}
    \item {{3}}
    \item {{8}}
\end{enumerate}
\vspace{2mm}

\noindent\textbf{Q60. Determine the volume of the solid enclosed by the cylinder x 2 + y 2 = 9 and the planes y + z = 5 and z = 1 .} \hfill [Tag: application of triple integrals to calculate volume | Hard]
\begin{enumerate}[label=(\alph*),topsep=2pt,itemsep=1pt,partopsep=0pt,parsep=0pt]
    \item {{36 $\pi$}}
    \item {{18 $\pi$}}
    \item {{54 $\pi$}}
    \item {{72 $\pi$}}
\end{enumerate}
\vspace{2mm}


\section{Section B: Comprehensive Subjective Problem with Detailed Analytical Rubric}

\begin{subjectiveq}{6}{1}
Change the order of integration and evaluate $I = \int_0^a \int_{x/a}^{\sqrt{x/a}} (x^2 + y^2) dy dx$.
\end{subjectiveq}

\begin{solbox}
The given region is bounded by $y = x/a$ (straight line) and $y^2 = x/a \iff x = ay^2$ (parabola). The intersection points are $(0,0)$ and $(a,1)$. Reversing the order of integration yields $x$ running from $ay^2$ to $ay$, while $y$ ranges from $0$ to $1$: $I = \int_0^1 \int_{ay^2}^{ay} (x^2 + y^2) dx dy = \int_0^1 \left[ \frac{x^3}{3} + xy^2 \right]_{ay^2}^{ay} dy$. Integrating term by term across $y \in [0, 1]$ yields $I = \frac{8a^3}{105}$.
\end{solbox}

\begin{examinertip}
Step 1: Identify the foundational definition and state governing theorems.\\
Step 2: Clearly write each algebraic transformation and boundary substitution.\\
Step 3: State the final numerical/algebraic invariant clearly in a concluding box.
\end{examinertip}
